\section{Related Work}
\label{sec:related}

This survey of related work concentrates on binary session types and their
treatment of recursion. Many works in the area concentrate on different aspects
and treat recursion informally as an add-on. In particular, the early work of
\citet{DBLP:conf/parle/TakeuchiHK94,DBLP:conf/concur/Honda93} did not consider
recursive types.
%
The school of session types that takes its foundations from the
sequent calculus for dual intuitionistic logic
\cite{DBLP:conf/concur/CairesP10,DBLP:journals/mscs/CairesPT16}
initially focuses on finite behaviors. While
recursion would break the logical consistency of their system, they
consider dependently typed variants that include replication
\cite{DBLP:conf/cpp/PfenningCT11}. 

\smallskip
\emph{Session type systems with recursion.}
Generally, recursive types come in two flavors, equirecursive and
isorecursive \cite{DBLP:conf/lics/AbadiF96}, and the impact of
these different approaches on type checking is well-studied
\cite{DBLP:books/daglib/0005958}. They have been found to be equally 
expressive \cite{DBLP:journals/pacmpl/PatrignaniMD21}. Most
session type systems in the literature support equirecursive session
types restricted to tail recursion and treat them informally.

Recursive types have first been introduced to session types by
\citet{DBLP:conf/esop/HondaVK98} in the context of $\pi$-calculus. Their
system relies on equirecursion and is syntactically restricted to
tail recursion.
\citet{DBLP:conf/esop/GayH99,DBLP:journals/acta/GayH05} study
subtyping in the same setting using a coinductive definition. 
% They
% define an algorithm for subtyping and develop a sound type checker.
% Most subsequent work refers to this treatment and the algorithm
% extends to their setting.
%
\citet{DBLP:conf/ppdp/CastagnaDGP09} consider recursive session types
as infinite regular trees. Their types are tagless in the sense that
choices are made according to the set of accepted values and the viability
of the continuation.
%
The functional session type calculus of
\citet{DBLP:journals/jfp/GayV10} includes subtyping and equirecursive
types, but the latter are restricted to tail-recursive session types.
%
Fundamentals of session types \cite{DBLP:journals/iandc/Vasconcelos12}
proposes a range of $\pi$-calculus-based systems with increasing
complexity. One particular focus is the treatment of replication and unrestricted
channel ends, which requires recursive types. In this work, recursive
types are modeled by infinite regular trees, represented using the
familiar $\mu$ notation with contractive, tail recursive bodies.
% no subtyping

\citet{DBLP:conf/esop/ToninhoCP13} extend the linear-logic based line
of work and focus on the combination of functions and processes
including recursive processes and consequently recursive
types for processes and functions. The examples indicate the equirecursive variant.
Apparently, there is no restriction in the use of recursion and
sequencing is implemented using the tensor $\otimes$.
% Saving some space; also I do not quite understand the restriction
% However, linearity stops them to write a type
% like $\alpha\otimes\alpha$ when $\alpha$ is a session type:
% operationally, it would mean that a process spawns a copy of itself,
% which is not possible in their system (unless via replication, which
% is modeled with a different logical connective).
%
\citet{DBLP:conf/tgc/ToninhoCP14} cover a different aspect of
equirecursion. They consider corecursive protocol definitions and
analyze their productivity. The main question they pursue is when a
corecursive protocol produces the next value after finitely many
internal steps.
%
SILL \cite{DBLP:conf/fossacs/PfenningG15} is a process-based language design that
integrates linear and affine types as well as synchronous and
asynchronous communication. It builds directly on the linear logic interpretation of session
types. Its examples make use of polymorphism and equirecursive types in the form of tail
recursive type equations, neither of which are reflected in the type
theoretical presentation.
% Subsequent work on manifest sharing also
% takes polymorphism and equirecursion for granted \cite{DBLP:journals/pacmpl/BalzerP17}.
%
\citet{DBLP:conf/concur/BartolettiSZ14} offer a semantic analysis of
session types on the basis of labeled transition systems. Their
system features equirecursion in tail position.


% \begin{itemize}
% \item X equi, regular, pi-calculus \cite{DBLP:conf/esop/HondaVK98},
%   first with recursion?
% \item X equi, regular, pi-calculus, subtyping \cite{DBLP:conf/esop/GayH99}
% \item X equi, regular \cite{DBLP:journals/jfp/GayV10}
% \item X equi, regular, pi-calculus ``fundamentals''
%   \cite{DBLP:journals/iandc/Vasconcelos12}
% \item X equi, regular, process and functional
%   \cite{DBLP:conf/esop/ToninhoCP13}
% \item X equiv, regular, process \cite{DBLP:conf/tgc/ToninhoCP14}: this
%   is really corecursion, focuses on productivity where a process is
%   guaranteed to take finitely many internal steps and then engages in
%   the next communication
% \item equi, regular, pi-calculus \cite{DBLP:journals/corr/Dardha14}
%   (focus on encoding of recursive sessions in recursive linear pi
%   calculus, duality and recursion)
% \item X equi, regular, lts \cite{DBLP:conf/concur/BartolettiSZ14}
% \item X SILL: implicit use of regular equirecursion \cite{DBLP:conf/fossacs/PfenningG15}
% \item X iso(?), regular \cite{DBLP:conf/icfp/LindleyM16}
% \item X equi, regular \cite{DBLP:journals/pacmpl/BalzerP17}
% \item equi, regular \cite{lindley17:_light_funct_session_types}
% \item Padovani's FUSE with OCaml implementation. 
% \item equi, regular \cite{DBLP:conf/esop/VasconcelosCAM20}
% \item infinite regular trees \cite{DBLP:journals/pacmpl/CicconeP22}
% \item X defined by different kinds of systems of equations \cite{DBLP:conf/fossacs/GayPV22}
% \end{itemize}


\citet{DBLP:conf/icfp/LindleyM16} define two
functional session type systems with recursion, $\mu$GV and $\mu$CP, and show that
they are related by typing- and semantics-preserving translations. So
we only compare to $\mu$GV in the following. $\mu$GV is based on
a simply-typed linear lambda calculus extended with notions of recursion and
corecursion. These are defined as least and greatest (respectively) fixed points of
strictly positive functors --- analogous to a logically
meaningful version of algebraic (co-) datatypes --- with generic fold and
unfold operators. Finally, they lift this isorecursive treatment of
recursive data to the level of session types by encoding branching and
recursion, mapping the corresponding data functors over session
types and wrapping them in promises.
We conclude that $\mu$GV features session types with isorecursion and
that it could be used to encode the simply-typed fragment of \algst,
but without parameterized data types and without the over
direction. Where $\mu$GV uses promises (i.e., a separate channel) to encode the argument to the
fold of the recursive type (i.e., the constructor arguments of an
algebraic datatype), \algst transmits the body of the fold on the same
channel, thus avoiding the overhead of channel creation.
Moreover, $\mu$GV is mainly a theoretical contribution
whereas a full implementation of \algst is available.

\smallskip
\emph{Context-free and nested session type systems.}
\citet{DBLP:conf/icfp/ThiemannV16} proposed
context-free session types to enable the type-safe encoding of
non-regular protocols like the serialization of tree structures or the
stack protocol from the introduction without encurring the overhead of
session initiation. They changed the asymmetric structure of the
traditional session type syntax $\TIn TS$ and $\TOut TS$, where an
item of type $\typec T$ is communicated and the session continues with
$\typec S$,
to a symmetric structure where session types can be built from the
atomic transmissions, $\TGIn T$ and $\TGOut T$, choice and branching,
$\TChoice{S_1}{S_2}$ and $\TBranch{S_1}{S_2}$,  and a monoidal sequential
composition operator $\TSeq{S_1}{S_2}$ with unit $\TSkip$. 
% For example, here
% is the type for transmitting a binary tree
% (\cf \cref{sec:algebr-datatyp-as}) in their system:
% $ \typec{\mu X.} \TChoice{\TGOut\TInt}{(\TSeq XX)}$.
This design,
together with equirecursive sessions, leads to a number of technical
problems. First, types come with a non-trivial equational theory:
sequential composition forms a monoid with unit $\TSkip$ and it
distributes over choice and branching. Second, polymorphic recursion
is needed as a recursive call may have to be instantiated with
different continuation sessions. Third, the interaction between
equirecursion and the omission of the restriction to tail recursion
gives rise to a difficult, but decidable, type equivalence relation.
%
Subsequent work investigated tractable algorithms for type
equivalence
\cite{DBLP:journals/corr/abs-1904-01284,DBLP:conf/tacas/AlmeidaMV20}.
\citet{DBLP:journals/corr/abs-2106-06658} integrate context-free
session types with explicit polymorphism in the style of
System~F. Recent work extends this approach with higher-order types
\cite{DBLP:journals/corr/abs-2203-12877}.

\citet{DBLP:conf/esop/Padovani17,DBLP:journals/toplas/Padovani19}
proposes a different approach to context-free session types embodied
in a functional session calculus called \fuse. This calculus uses the
type language of context-free session types, but
sidesteps the problems with the sequencing operator and type
equivalence by introducing a resumption operator $\{e\}_c$ that
deconstructs a sequence of sessions on channel $c :
\TSeq{S_1}{S_2}$. It does so by retyping the channel to $c : \typec{S_1}$ for
evaluating $e$, which must return the same channel after finishing
$\typec{S_1}$.
More formally, the resumption obeys some kind of frame rule: if  $c :
\typec{S_1} \vdash e : \typec{T \times \TOne}$, then
$c : \TSeq{S_1}{S_2} \vdash \{e\}_c : \typec{T \times S_2}$. The
resumption temporarily removes the continuation
$\typec{S_2}$ from the session type and sticks it back on
afterwards. The type system has means to make sure that the channel returned by $e$ is
identical to $c$, which is required for soundness. The type $\TOne$
does not signify the end of the session, but rather that some known
part of the session has been processed and the session will be
continued in the context (of the resumption).
%
The management of the choice and branch operations is very close to
\algst. \fuse implements choice by transmitting the constructor tag in
a variant type that selects the different continuations. Its branch
operator matches on the received constructor and thus selects the
desired continuation. However, the handling of protocol constructors
in \algst includes the construction of the continuation type from the
types and directions given in the protocol declaration. \fuse has no
equivalent to a protocol declaration nor to the type operator
\lstinline+-+ to swap the direction of a transmission.
%
\citet{DBLP:journals/toplas/Padovani19} implements
\fuse in OCaml and demonstrates that context-free session types
are amenable to type inference, if the programmer helps with judicious
placement of resumptions. Support for equirecursion is already built into  OCaml's type inference
engine. It is also shown that subtyping is undecidable.

\citet{DBLP:conf/esop/DasDMP21} propose the metatheory of 
nested session types featuring parametrized definitions of session types
with nested polymorphic recursion. 
At first glance, the nested session type 
$\syntax{List}[\alpha] = \oplus\{\syntax{nil}\colon \syntax{end}, \syntax{cons}\colon \alpha \otimes \syntax{List}[\alpha]\}$
is very similar to the algebraic protocol
$\protocolk\ \TPApp{\syntax{List}}\alpha = \termc{\syntax{Nil}}\ |\
\termc{\syntax{Cons}}\ \alpha\ (\typec{\syntax{List}}\ \alpha)$, but
there are significant differences. On the
algebraic side, we can materialize the \typec{\syntax{List}} protocol in a type  
that sends or receives a list of integers, while the nested type $\syntax{List}$ is 
already committed to a concrete direction.
%
Although type nesting enables to capture the modularity characterizing 
algebraic sessions, there are several features that distinguish AlgST from 
nested session types:
Das et al. propose a structural and equirecursive, rather than a
nominal, view of type definitions, and interpret types coinductively, rather than inductively.
These features have significant impact on the complexity of the 
equivalence problem for nested session types and on its implementation:
Das et al. present a sound but incomplete algorithm for type equivalence.
By considering a nominal interpretation of algebraic session types, we reduced the 
complexity of the equivalence problem from 
doubly-exponential~\cite{DBLP:conf/esop/DasDMP21} to 
linear and designed a sound and complete algorithm to verify type
equivalence. The examples in section~2 of
\citet{DBLP:conf/esop/DasDMP21} translate into \algst and type check with our implementation.

\citet{DBLP:conf/fossacs/GayPV22} consider equirecursive protocols
defined by different kinds of systems of equations (with and without
parameters). Imposing different restrictions on the parameterization
yields different classes of protocols with decision problems for type
equivalence ranging from polytime to undecidable. Their largest
decidable class of pushdown session types is shown to be equivalent to
nested session types, which we can model with \algst.

Compared to all competiting systems \cite{DBLP:journals/pacmpl/ThiemannV20,DBLP:journals/corr/abs-2106-06658,DBLP:journals/toplas/Padovani19,DBLP:conf/esop/DasDMP21}, \algst's materialization and
directional operators are unique. These
operators enable the construction of protocols without committing to a
particular direction. 

%%% MPST:
%%% (equi, regular) Less is more: multiparty session types revisited

% \paragraph{Protocol composition}
% \label{sec:protocol-composition}

% \citet{DBLP:journals/corr/abs-2203-02461} define a notion of protocol
% composition which interleaves two or more protocols. To
% restrict this liberal approach to sensible interleavings, protocols
% are extended with an assertion language that specifies safe
% interleavings. Interleaving composition is complementary to our
% modular approach to protocol construction: we combine subprotocols
% rather than interleave them.

\smallskip
\emph{Linear and unrestricted types.}
Our 
%formal 
calculus is mostly linear, except that we
%. The only exception we make is to
allow unrestricted bindings for recursive functions, which would be
rendered useless otherwise. 
%On the other hand, we want to speak about
%recursive protocols, so we felt like recursive functions are needed in
%the formal development. 
Our implementation goes well beyond that in
allowing fully fledged unrestricted computations besides ensuring that
session typed channels are treated linearly.
% That said, all related work is closer to our implementation than to
% the formal calculus.
%
There is some work on recursion and linearity in the context of lambda calculus
\cite{DBLP:journals/lisp/AlvesFFM10}
%. This work 
adding iterators on
linear natural numbers to the calculus, thereby avoiding the need for
a rule like {\rulenameexpUVar}.
%, which we need to bind the recursive function to a name.

System F\degree~\cite{DBLP:conf/tldi/MazurakZZ10}
extends System F with kinds to integrate polymorphism with linear and
unrestricted types. This work serves as a blueprint for much
subsequent work in this area including LFST and FreeST (see below), as well as our
implemented language.
They categorize all objects, including functions, into one-use or
many-use resources. This style can be traced back to 
\citet{walker:substructural-type-systems}. 
% One example in the paper applies the language to protocols built with
% linear capabilities, reminiscent of session types.
%
LFST \cite{lindley17:_light_funct_session_types} formalizes the
session type system of Links, which features a mix of linear and
unrestricted types, embedding a standard functional programming
language. The substructural behavior of types is specified by a
sophisticated kind structure.  The kinding used in our implementation
is inspired by LFST, though LFST has additional kinds to describe row
types that we do not support.
%
\citet{DBLP:journals/corr/abs-2106-06658}
consider an integration of context-free session types with
System~F. Their type system includes a kinding discipline that manages
linear and unrestricted versions of ordinary types and
message types (possible payloads for the send and receive
operations). Our implemented system extends the kind structure in a
similar way.

Linear Haskell (LH) \cite{DBLP:journals/pacmpl/BernardyBNJS18}
is a proposal to integrate linear types with stock functional
programming. While the work discussed so far characterizes resources,
LH supports the lollipop type $S \multimap T$ which classifies
functions that use their argument exactly once. Originally, LH forced
the programmer to allocate resources using continuations. This restriction
has been fixed in subsequent work \cite{DBLP:journals/pacmpl/SpiwackKBWE22}.
%
There is further work on integrating linear and dependent types (e.g.,
\cite{DBLP:conf/popl/KrishnaswamiPB15}) which relies on
bifurcating the language in two calculi connected by an adjunction.

% We believe that \algst definitions that do not make
% significant use of negation on protocols

%%% Local Variables:
%%% mode: latex
%%% TeX-master: "main"
%%% End:


% \paragraph{Further work on duality and recursion}

% The problem with duality and recursion mentioned in
% \cref{sec:what-about-inter} was first reported by
% \citet{DBLP:journals/corr/abs-1202-2086}. % (Bono and Padovani, 2012)
% It was further studied and fixed by \citet{BernardiH14, BernardiH16}. %(Bernardi Hennessy, 2014)
% \citet{DBLP:conf/icfp/LindleyM16} propose a different solution using
% covariables. A recent study closes the matter \cite{DBLP:journals/corr/abs-2004-01322}.



% (Gay and Thiemann and Vasconcelos, 2020)
% \cite{DBLP:journals/corr/Dardha14}

% \paragraph{Not recursive}
% \begin{itemize}
% \item Towards Races in Linear Logic (Kokke, Morris, Wadler)
% \item Exceptional GV (Fowler etal)
% \item Gradual Session Types 
% \item Bernardo Toninho, Nobuko Yoshida: Interconnectability of
%   Session-Based Logical Processes
% \item Bernardo Toninho, Nobuko Yoshida: Depending on Session-Typed
%   Processes
% \item Sam Lindley, J. Garrett Morris: Embedding session types in Haskell.
% \item Lindley and Morris, Semantics for Propositions as Sessions, ESOP
%   2015
% \item Mazurak and Zdancewic, Lolliproc: to Concurrency from Classical
% Linear Logic via Curry-Howard and Control, ICFP 2010
% \item Perez and Caires and Pfenning and Toninho, Linear logical
%   relations and observational equivalences for session-based
%   concurrency, I\&C 2014
% \item Riccardo Pucella and Jesse A. Tov, Haskell Session Types with
%   (Almost) No Class, 2008
% \item X An interaction-based language and its typing system
% Kaku Takeuchi, Kohei Honda \& Makoto Kubo, PARLE 1994 \cite{DBLP:conf/parle/TakeuchiHK94}
% \end{itemize}

%%% Local Variables:
%%% mode: latex
%%% TeX-master: "main"
%%% End:
